\documentclass[11pt]{beamer}
\usetheme{Frankfurt}
\usecolortheme{crane}
\usepackage[utf8]{inputenc}
\usepackage[english]{babel}
\usepackage{amsmath}
\usepackage{amsfonts}
\usepackage{amssymb}
\usepackage{graphicx}
%\usepackage{circuitikz}
\author{Maximilian Heim}
\title{Hardware Trojan Detection}
%\setbeamercovered{transparent} 
%\setbeamertemplate{navigation symbols}{} 
\logo{\includegraphics[width=1.5cm]{1200px-Hsas_logo.svg.png}}
\institute{University Albstadt-Sigmaringen} 
\date{\today} 
\subject{Hardware Cyber-Security}
\begin{document}

\begin{frame}
\titlepage
\end{frame}

\begin{frame}
\tableofcontents
\end{frame}

\section{Introduction}
\subsection{Hardware Trojans}
\begin{frame}
    \frametitle{Was sind Hardware Trojaner?}
    \begin{enumerate}
    \item Bösartige Modifikation eines Integrierten Schaltkreises
    \item Besteht aus einem Trigger und einer Payload
    \item Triggers: Time bombs, cominational, analog \dots
    \item Payloads: Denial of service, extraction of information, keys
    \end{enumerate}
\end{frame}
\subsection{Relevanz}
\begin{frame}
    \frametitle{Warum ist es wichtig diese zu erkennen?}
    \begin{itemize}
    \item Militär
    \item Finanzen
    \item Energie
    \item Geheimdienste
    \item Überwachungsstaaten
    \item Transport
    \item U.v.m \ldots
    \end{itemize}

\end{frame}


\section{Erkennung}
\subsection{Destruktive Erkennung}
\begin{frame}
    \frametitle{Reverse Engineering}
    \begin{itemize}
        \item Entfernen der Oberfläche Schicht für Schicht
        \item Advantages: 
        \begin{enumerate}
            \item 100 \% detection rate
        \end{enumerate}
        \item Disadvantages:
        \begin{enumerate}
            \item Only tests a single chip
            \item Chip unusable afterwards
            \item Very time consuming
        \end{enumerate}
        \item However, 
    \end{itemize}
\end{frame}
\subsection{Nicht-Destruktive Erkennung}
\section{Conclusion}
\begin{frame}
    \frametitle{Types of detection}
\end{frame}

\end{document}